% ============================================================
%  Explicacion_pesos.tex
%  Explicación detallada: Tablas 7-13 de Comparativa
%  Pesos hat{omega} y hat{lambda} — Python vs Stata
% ============================================================
\documentclass[11pt, a4paper]{article}

\usepackage[utf8]{inputenc}
\usepackage[T1]{fontenc}
\usepackage[margin=2.5cm]{geometry}
\usepackage{booktabs}
\usepackage{graphicx}
\usepackage{float}
\usepackage{amsmath}
\usepackage{amssymb}
\usepackage{array}
\usepackage{xcolor}
\usepackage{colortbl}
\usepackage{multirow}
\usepackage{hyperref}
\usepackage{lmodern}
\usepackage{microtype}
\usepackage{parskip}
\usepackage{tcolorbox}
\usepackage{enumitem}
\usepackage{caption}
\usepackage[spanish]{babel}

\hypersetup{colorlinks=true, linkcolor=blue!60!black, urlcolor=blue!60!black}

\definecolor{pyblue}{RGB}{31,119,180}
\definecolor{stgreen}{RGB}{44,160,44}
\definecolor{headgray}{RGB}{235,235,235}
\definecolor{highlight}{RGB}{255,250,205}
\definecolor{sdidcolor}{RGB}{204,229,255}
\definecolor{sccolor}{RGB}{204,255,204}
\definecolor{didcolor}{RGB}{255,229,204}

\newcommand{\py}{\textcolor{pyblue}{\textbf{Python}}}
\newcommand{\st}{\textcolor{stgreen}{\textbf{Stata}}}
\newcommand{\sdid}{\textsc{sdid}}
\newcommand{\scc}{\textsc{sc}}
\newcommand{\did}{\textsc{did}}

\captionsetup{font=small, labelfont=bf}
\tcbuselibrary{skins, breakable}

% Caja de concepto
\newtcolorbox{concepto}[1]{
  colback=blue!5!white, colframe=blue!50!black,
  fonttitle=\bfseries, title=#1, breakable
}
% Caja de interpretación
\newtcolorbox{interpretacion}{
  colback=yellow!10!white, colframe=orange!60!black,
  fonttitle=\bfseries, title=Interpretación, breakable
}
% Caja de resultado
\newtcolorbox{resultado}{
  colback=green!8!white, colframe=green!50!black,
  fonttitle=\bfseries, title=Resultado clave, breakable
}

% ============================================================
\begin{document}

\begin{titlepage}
  \centering
  \vspace*{2.5cm}
  {\LARGE\bfseries Explicación de Pesos\\[0.5em]
  Tablas 7--13 del Documento Comparativa\\[1em]}
  {\large Synthetic Difference-in-Differences (SDID)\\
  vs Synthetic Control (SC) vs Difference-in-Differences (DiD)\\[1.5em]}
  {\normalsize
    \textbf{Dataset:} Cigar (Baltagi, 2002) — 46 estados EE.UU., 1963--1992\\[0.3em]
    \textbf{Outcome:} Ventas de cigarrillos (paquetes per cápita)\\[0.3em]
    \textbf{Comparación:} \py{} (\texttt{synthdid.py}) vs \st{} (\texttt{sdid})\\[2em]
  }
  \vfill
  {\small Marzo 2026}
\end{titlepage}

\tableofcontents
\newpage

% ============================================================
\section{Marco Conceptual: ¿Qué son los pesos $\hat{\omega}$ y $\hat{\lambda}$?}
% ============================================================

El estimador SDID (y sus variantes SC y DiD) construye un \textbf{contrafactual sintético}
combinando unidades de control y periodos de tiempo mediante dos tipos de pesos:

\begin{concepto}{Pesos de unidades de control $\hat{\omega}_i$ (quién)}
Determinan \textbf{qué estados de control} contribuyen al contrafactual y en qué proporción.
Un $\hat{\omega}_i$ alto significa que el estado $i$ es muy parecido al grupo tratado
en el período pre-tratamiento. Se exige que $\hat{\omega}_i \geq 0$ y $\sum_i \hat{\omega}_i = 1$.
\end{concepto}

\begin{concepto}{Pesos temporales $\hat{\lambda}_t$ (cuándo)}
Determinan \textbf{qué años pre-tratamiento} son más informativos para predecir el
nivel contrafactual justo antes del tratamiento. Se exige que $\hat{\lambda}_t \geq 0$
y $\sum_t \hat{\lambda}_t = 1$. Solo existen en SDID y DiD (en SC son todos cero).
\end{concepto}

Cada método los obtiene de forma diferente:

\begin{center}
\begin{tabular}{lccc}
\toprule
& \textbf{SDID} & \textbf{SC} & \textbf{DiD} \\
\midrule
$\hat{\omega}_i$ & Optimización regularizada & Optimización (sin regularizar) & Uniforme $1/N_0$ \\
$\hat{\lambda}_t$ & Optimización regularizada & $0$ para todo $t$ & Uniforme $1/T_0$ \\
\bottomrule
\end{tabular}
\end{center}

En el dataset Cigar hay \textbf{38 estados de control} ($N_0=38$) y \textbf{7 cohortes}
de adopción: 1975, 1976, 1978, 1979, 1981, 1983 y 1986.
Cada cohorte tiene su propia estimación de $\hat{\omega}$ y $\hat{\lambda}$ de forma independiente.

% ============================================================
\newpage
\section{Tabla 7 — Diferencia máxima entre pesos Python y Stata}
% ============================================================

\begin{table}[H]
\centering
\caption{Diferencia máxima absoluta entre pesos \py{} y \st{}}
\begin{tabular}{lcc}
\toprule
\textbf{Método} & $\max|\hat{\omega}^{\text{Py}} - \hat{\omega}^{\text{Stata}}|$ &
                  $\max|\hat{\lambda}^{\text{Py}} - \hat{\lambda}^{\text{Stata}}|$ \\
\midrule
\rowcolor{sdidcolor}
SDID & $4.8 \times 10^{-8}$ & $4.0 \times 10^{-8}$ \\
\rowcolor{sccolor}
SC   & $5.0 \times 10^{-8}$ & $0$ (exacto) \\
\rowcolor{didcolor}
DiD  & $1.1 \times 10^{-8}$ & $4.4 \times 10^{-8}$ \\
\bottomrule
\end{tabular}
\end{table}

\subsection*{¿Qué mide esta tabla?}

Esta tabla es el \textbf{resumen ejecutivo} de toda la comparación de pesos.
Para cada método, reporta la diferencia absoluta máxima entre los pesos calculados
por \py{} y los calculados por \st{}, tomando el máximo sobre todas las cohortes
y todas las unidades (o años).

\begin{interpretacion}
Las diferencias son del orden de $10^{-8}$, es decir, $0.00000005$. Esto es
\textbf{ruido numérico puro}, no una diferencia real de implementación.

¿Por qué existen estas diferencias? Ambos paquetes usan el algoritmo
Frank-Wolfe para resolver la optimización, pero lo implementan en entornos distintos:
\begin{itemize}
  \item \py{}: NumPy (Python), aritmética en C/Fortran con IEEE 754 double precision
  \item \st{}: Mata (Stata), su propio motor de álgebra matricial
\end{itemize}
Pequeñas diferencias en el orden de las operaciones flotantes se acumulan iteración
a iteración hasta producir discrepancias de $\sim10^{-8}$.
\end{interpretacion}

\begin{resultado}
A \textbf{4 decimales}, los pesos son \textbf{idénticos} entre Python y Stata
en los tres métodos y las siete cohortes. El máximo diferencial ($5 \times 10^{-8}$)
equivale a un error de 1 en la octava cifra decimal. Ambas implementaciones
convergen al mismo óptimo matemático.

Nota especial SC: $\hat{\lambda}^{SC} = 0$ exactamente en ambos paquetes
porque es un valor fijo por diseño, no un resultado numérico.
\end{resultado}

% ============================================================
\newpage
\section{Tablas 8, 9 y 10 — Pesos temporales $\hat{\lambda}$ por método}
% ============================================================

\subsection{Tabla 8 — $\hat{\lambda}$ SDID}

\begin{table}[H]
\centering
\caption{Pesos temporales $\hat{\lambda}$ --- SDID. Python $=$ Stata a 4 decimales.}
\small
\begin{tabular}{rrrrrrrr}
\toprule
Año & 1975 & 1976 & 1978 & 1979 & 1981 & 1983 & 1986 \\
\midrule
1963 & 0       & 0.0011 & 0.2195 & 0.1566 & 0.0641 & 0.0711 & 0      \\
1965 & 0.1382  & 0.1854 & 0      & 0      & 0      & 0      & 0      \\
1969 & 0       & 0.0034 & 0      & 0      & 0      & 0      & 0      \\
1974 & \cellcolor{highlight}\textbf{0.8618}  & 0      & 0      & 0      & 0      & 0      & 0      \\
1975 & --      & \cellcolor{highlight}\textbf{0.8100} & 0      & 0      & 0      & 0      & 0      \\
1977 & --      & --     & \cellcolor{highlight}\textbf{0.7805} & 0      & 0      & 0      & 0      \\
1978 & --      & --     & --     & \cellcolor{highlight}\textbf{0.8434} & 0      & 0      & 0      \\
1980 & --      & --     & --     & --     & \cellcolor{highlight}\textbf{0.9359} & 0      & 0      \\
1982 & --      & --     & --     & --     & --     & \cellcolor{highlight}\textbf{0.9289} & 0      \\
1985 & --      & --     & --     & --     & --     & --     & \cellcolor{highlight}\textbf{1.0000} \\
\midrule
\multicolumn{8}{l}{\textit{Los años no mostrados tienen $\hat{\lambda}_t = 0$.
``--'' indica año post-tratamiento para esa cohorte.}} \\
\bottomrule
\end{tabular}
\end{table}

\subsubsection*{¿Qué significa?}

El SDID asigna más peso a los años pre-tratamiento que mejor predicen el nivel
de la variable de resultado justo \textit{antes} de la intervención. El patrón
es muy claro:

\begin{itemize}
  \item \textbf{El año inmediatamente anterior al tratamiento recibe casi todo el peso}
    (resaltado en amarillo), entre 0.76 y 1.00.
  \item Los demás años reciben pesos pequeños o cero.
\end{itemize}

\begin{interpretacion}
\textbf{¿Por qué el último año pre-tratamiento domina?}

El algoritmo SDID busca los años donde el control sintético (construido con
$\hat{\omega}$) se parece más al grupo tratado. El año más reciente antes del
tratamiento tiende a ser el mejor predictor del nivel contrafactual
``en ausencia del tratamiento'', porque captura las tendencias de corto plazo.

\medskip
\textbf{Cohorte 1986} es el caso extremo: $\hat{\lambda}_{1985}=1.0$.
Todo el peso recae en el año 1985, el único que predice perfectamente el nivel
de ventas de cigarrillos que habría tenido el estado 42 si no hubiera subido
el impuesto en 1986.

\medskip
\textbf{Cohorte 1963 vs 1965:} Algunas cohortes (1976, 1978, 1979...) tienen
pesos residuales en 1963 o 1965. Esto ocurre cuando hubo una tendencia especial
en esos años que el algoritmo considera relevante para el matching de largo plazo.

\medskip
En los \textbf{gráficos de trayectorias}, el área verde sombreada
\textit{es} este $\hat{\lambda}$: las barras más altas corresponden a los años
con más peso, mostrando visualmente qué parte del pre-período ``importa'' más.
\end{interpretacion}

\subsection{Tabla 9 — $\hat{\lambda}$ SC}

\begin{table}[H]
\centering
\caption{Pesos temporales $\hat{\lambda}$ --- SC. Python $=$ Stata a 4 decimales.}
\small
\begin{tabular}{rrrrrrrr}
\toprule
Año & 1975 & 1976 & 1978 & 1979 & 1981 & 1983 & 1986 \\
\midrule
\multicolumn{8}{c}{\textit{SC: $\hat{\lambda}_t = 0\ \forall\,t$ (sin ponderación temporal pre-período).}} \\
\bottomrule
\end{tabular}
\end{table}

\subsubsection*{¿Qué significa?}

El Synthetic Control puro \textbf{no usa pesos temporales}. Fija $\hat{\lambda}_t = 0$
para todos los años pre-tratamiento por diseño. El matching se hace únicamente
a través de los pesos de unidades $\hat{\omega}_i$.

\begin{interpretacion}
En el SC clásico (Abadie et al., 2010), el objetivo es encontrar una combinación
de controles que replique el \textit{nivel promedio} del outcome pre-tratamiento
(y a veces covariables adicionales). No existe una selección de qué años importan
más: todos los periodos pre-tratamiento cuentan por igual \textit{en la restricción
de matching}, aunque el promedio no se construye con pesos lambda explícitos.

\medskip
En términos de los gráficos de trayectorias: la cohorte SC
\textbf{no muestra área verde sombreada} porque $\hat{\lambda}_t = 0\ \forall\,t$.
\end{interpretacion}

\subsection{Tabla 10 — $\hat{\lambda}$ DiD}

\begin{table}[H]
\centering
\caption{Pesos temporales $\hat{\lambda}$ --- DiD. Python $=$ Stata a 4 decimales.}
\small
\begin{tabular}{rrrrrrrr}
\toprule
Año & 1975 & 1976 & 1978 & 1979 & 1981 & 1983 & 1986 \\
\midrule
1963 & 0.0833 & 0.0769 & 0.0667 & 0.0625 & 0.0556 & 0.0500 & 0.0435 \\
1964 & 0.0833 & 0.0769 & 0.0667 & 0.0625 & 0.0556 & 0.0500 & 0.0435 \\
$\vdots$ & $\vdots$ & $\vdots$ & $\vdots$ & $\vdots$ & $\vdots$ & $\vdots$ & $\vdots$ \\
1974 & 0.0833 & 0.0769 & 0.0667 & 0.0625 & 0.0556 & 0.0500 & 0.0435 \\
1975 & --     & 0.0769 & 0.0667 & 0.0625 & 0.0556 & 0.0500 & 0.0435 \\
$\vdots$ & & $\vdots$ & $\vdots$ & $\vdots$ & $\vdots$ & $\vdots$ & $\vdots$ \\
1985 & --     & --     & --     & --     & --     & --     & 0.0435 \\
\midrule
$T_0$ & 12 & 13 & 15 & 16 & 18 & 20 & 23 \\
$1/T_0$ & 0.0833 & 0.0769 & 0.0667 & 0.0625 & 0.0556 & 0.0500 & 0.0435 \\
\bottomrule
\end{tabular}
\end{table}

\subsubsection*{¿Qué significa?}

El DiD asigna \textbf{peso uniforme} $\hat{\lambda}_t = 1/T_0$ a cada año
pre-tratamiento. Cada cohorte tiene un valor distinto porque tiene distinto $T_0$:

\begin{itemize}
  \item Cohorte 1975: $T_0=12$ años pre $\Rightarrow$ $\hat{\lambda}_t = 1/12 \approx 0.0833$
  \item Cohorte 1986: $T_0=23$ años pre $\Rightarrow$ $\hat{\lambda}_t = 1/23 \approx 0.0435$
\end{itemize}

\begin{interpretacion}
El DiD \textbf{trata todos los años pasados como igualmente informativos}.
Esto equivale a comparar diferencias en promedio pre vs promedio post,
sin distinguir si 1963 o 1974 es más relevante para predecir 1975.

Esta restricción es más fuerte que SDID (que concentra el peso en 1-2 años)
y hace al DiD más sensible a tendencias pre-tratamiento no paralelas.
\end{interpretacion}

% ============================================================
\newpage
\section{Tablas 11, 12 y 13 — Pesos de unidades $\hat{\omega}$ por método}
% ============================================================

\subsection{Tabla 11 — $\hat{\omega}$ SDID}

\begin{table}[H]
\centering
\caption{Pesos unitarios $\hat{\omega}$ --- SDID. Python $=$ Stata a 4 decimales. (Selección de estados)}
\small
\begin{tabular}{rrrrrrrr}
\toprule
Estado & 1975 & 1976 & 1978 & 1979 & 1981 & 1983 & 1986 \\
\midrule
1  & 0.0385 & 0      & 0.0463 & 0.0084 & 0      & 0.0131 & 0.0352 \\
5  & 0.0503 & \textbf{0.1436} & 0.0078 & 0      & 0.0979 & 0.0173 & 0.0082 \\
8  & 0      & 0.0126 & 0      & 0.0291 & \textbf{0.2680} & 0.0520 & 0.0174 \\
14 & 0.0387 & \textbf{0.1770} & 0.0041 & 0.0038 & \textbf{0.1180} & 0.0284 & 0.0308 \\
18 & 0.0579 & 0      & \textbf{0.1442} & 0.0326 & 0      & 0.0062 & 0.0257 \\
22 & 0.0296 & 0.0637 & 0.0035 & 0.0021 & \textbf{0.2476} & 0.0375 & 0.0238 \\
31 & 0.0137 & \textbf{0.2523} & 0      & 0.0102 & \textbf{0.1750} & 0.0150 & 0.0134 \\
36 & 0      & \textbf{0.1110} & 0      & 0.0298 & 0.0299 & 0.0527 & 0.0196 \\
39 & 0.0305 & \textbf{0.1293} & 0.0095 & 0.0265 & 0.0068 & 0.0378 & 0.0464 \\
46 & 0.0535 & 0      & 0.0752 & 0.0395 & 0      & 0.0085 & 0.0718 \\
$\vdots$ & $\vdots$ & $\vdots$ & $\vdots$ & $\vdots$ & $\vdots$ & $\vdots$ & $\vdots$ \\
\midrule
\multicolumn{8}{l}{\textit{En negrita: pesos $\geq 0.10$. Total: 38 estados de control, todos con $\hat{\omega}_i\geq 0$, $\sum=1$.}} \\
\bottomrule
\end{tabular}
\end{table}

\subsubsection*{¿Qué significa?}

Los pesos $\hat{\omega}^{SDID}$ miden la importancia de cada estado de control
en la construcción del contrafactual sintético. Sus características son:

\begin{enumerate}
  \item \textbf{Difusos}: La mayoría de los 38 estados recibe algún peso positivo.
    Raramente un solo estado domina.
  \item \textbf{Regularizados}: El SDID penaliza pesos extremos mediante el parámetro
    $\zeta_\omega = \eta \cdot \hat{\sigma}$, donde $\eta = (N_1 T_1)^{1/4}$.
    Esto evita el overfitting al pre-período.
  \item \textbf{Varían por cohorte}: Cada cohorte tiene su propio grupo de controles
    más relevantes, porque el conjunto de tratados cambia.
\end{enumerate}

\begin{interpretacion}
\textbf{Cohorte 1976} es notable: el estado 31 recibe $\hat{\omega}_{31}=0.2523$
(el peso más alto del SDID en todo el documento). Esto indica que ese estado
de control sigue una trayectoria de ventas de cigarrillos muy similar a los
2 estados tratados en 1976 (estados 9 y 40).

\medskip
\textbf{Cohorte 1981}: Los estados 8 (0.2680) y 22 (0.2476) dominan,
concentrando juntos el 51.6\% del peso. Esto sucede porque, para esa cohorte,
esos dos estados son los mejores ``gemelos'' del estado tratado (7)
en términos de tendencia pre-tratamiento de ventas.

\medskip
Los estados con $\hat{\omega}_i=0$ no contribuyen al contrafactual;
el algoritmo los descartó porque sus trayectorias no aportan información
adicional para el matching.
\end{interpretacion}

\subsection{Tabla 12 — $\hat{\omega}$ SC}

\begin{table}[H]
\centering
\caption{Pesos unitarios $\hat{\omega}$ --- SC. Python $=$ Stata a 4 decimales. (Selección de estados)}
\small
\begin{tabular}{rrrrrrrr}
\toprule
Estado & 1975 & 1976 & 1978 & 1979 & 1981 & 1983 & 1986 \\
\midrule
4  & 0      & 0      & 0      & \textbf{0.4494} & 0      & 0      & 0.0162 \\
5  & 0.1193 & 0      & 0      & 0      & 0      & 0      & 0      \\
8  & 0      & 0.1676 & 0.0030 & 0      & 0.0435 & 0.0285 & 0      \\
10 & 0.0302 & 0      & 0.0885 & 0      & 0      & 0.1354 & 0      \\
14 & 0      & \textbf{0.3841} & 0      & 0      & 0      & 0      & 0      \\
17 & \textbf{0.3106} & 0      & 0      & 0.0129 & 0      & 0      & 0      \\
18 & 0.0480 & 0      & \textbf{0.3175} & 0.0156 & 0      & 0.0147 & 0.0002 \\
22 & 0      & 0      & 0      & 0      & \textbf{0.9188} & 0.0026 & 0      \\
30 & 0.0005 & \textbf{0.3057} & 0.0335 & 0.0448 & 0      & 0.0562 & 0.0036 \\
33 & 0.0882 & 0      & \textbf{0.3645} & 0      & 0      & 0      & 0.0875 \\
45 & 0.0001 & 0.0007 & 0.0077 & 0.0083 & 0      & 0.0051 & \textbf{0.4037} \\
46 & 0.1933 & 0      & 0.1025 & 0.0474 & 0      & 0      & \textbf{0.2015} \\
50 & 0      & 0      & 0      & \textbf{0.2647} & 0      & \textbf{0.3445} & 0      \\
$\vdots$ & $\vdots$ & $\vdots$ & $\vdots$ & $\vdots$ & $\vdots$ & $\vdots$ & $\vdots$ \\
\midrule
\multicolumn{8}{l}{\textit{En negrita: pesos $\geq 0.10$. La mayoría de estados tienen $\hat{\omega}_i=0$.}} \\
\bottomrule
\end{tabular}
\end{table}

\subsubsection*{¿Qué significa?}

Los pesos $\hat{\omega}^{SC}$ son mucho más \textbf{dispersos} (sparse) que los del SDID:

\begin{enumerate}
  \item \textbf{Alta concentración}: 1-4 estados reciben casi todo el peso por cohorte.
  \item \textbf{Muchos ceros}: La mayoría de los 38 estados de control reciben $\hat{\omega}_i=0$.
  \item \textbf{Sin regularización} (o mínima, $\zeta=10^{-6}\hat{\sigma}$): el SC
    puede asignar pesos muy extremos si un estado replica bien el pre-período.
\end{enumerate}

\begin{interpretacion}
\textbf{Casos extremos:}

\begin{itemize}
  \item \textbf{Cohorte 1981}: Estado 22 recibe $\hat{\omega}_{22}^{SC}=0.9188$.
    Un solo estado de control construye el 92\% del contrafactual. Esto indica que
    el estado 22 tiene una trayectoria pre-tratamiento \textit{casi idéntica}
    al estado tratado 7, pero también es señal de potencial \textbf{overfitting}:
    el contrafactual ``memoriza'' el pre-período en lugar de generalizar.
  \item \textbf{Cohorte 1979}: Estado 4 recibe 0.4494, Estado 50 recibe 0.2647.
    Dos estados concentran el 71\% del peso.
\end{itemize}

\medskip
\textbf{Comparación con SDID:} El SDID distribuye el peso entre muchos más estados,
lo que le da mayor \textbf{estabilidad} ante posibles outliers. El SC puede dar
estimaciones más precisas si el matching es perfecto, pero más sensibles si uno
de los estados ``dominantes'' tiene un shock idiosincrásico.
\end{interpretacion}

\subsection{Tabla 13 — $\hat{\omega}$ DiD}

\begin{table}[H]
\centering
\caption{Pesos unitarios $\hat{\omega}$ --- DiD. Python $=$ Stata a 4 decimales.}
\small
\begin{tabular}{rrrrrrrr}
\toprule
Estado & 1975 & 1976 & 1978 & 1979 & 1981 & 1983 & 1986 \\
\midrule
\multicolumn{8}{c}{\textit{DiD: los 38 estados de control reciben peso uniforme
$\hat{\omega}_i = 1/38 \approx 0.0263$ en todas las cohortes.}} \\
\bottomrule
\end{tabular}
\end{table}

\subsubsection*{¿Qué significa?}

El DiD asigna el mismo peso $1/38 \approx 0.0263$ a \textbf{todos} los estados de control,
independientemente de si su trayectoria pre-tratamiento se parece o no a la del tratado.

\begin{interpretacion}
\textbf{El DiD es el estimador menos flexible de los tres.} No intenta encontrar
los controles más parecidos: simplemente promedia a todos por igual.
Esto equivale a comparar la diferencia en medias:

\[
  \widehat{ATT}^{DiD} = \underbrace{\bar{Y}_{\text{tratados, post}}}_{\text{promedio post trat.}}
  - \underbrace{\bar{Y}_{\text{tratados, pre}}}_{\text{promedio pre trat.}}
  - \left(\underbrace{\bar{Y}_{\text{controles, post}}}_{\text{promedio post control}}
         - \underbrace{\bar{Y}_{\text{controles, pre}}}_{\text{promedio pre control}}\right)
\]

\medskip
En el documento, el ATT DiD es $-19.48$ vs SDID $-13.85$.
La diferencia ($\approx 5.6$ paquetes) refleja que el grupo de controles
promedio \textit{no} es un buen contrafactual para los estados tratados:
los estados con alza fiscal tenían tendencias pre-tratamiento distintas
al promedio nacional, lo que el SDID y SC corrigen mediante sus pesos óptimos.
\end{interpretacion}

% ============================================================
\newpage
\section{Comparación cruzada de los tres métodos}
% ============================================================

\subsection{Resumen de pesos $\hat{\lambda}$}

\begin{table}[H]
\centering
\caption{Características de $\hat{\lambda}$ por método}
\begin{tabular}{lccc}
\toprule
\textbf{Característica} & \textbf{SDID} & \textbf{SC} & \textbf{DiD} \\
\midrule
Optimización & Sí (regularizada) & No (fijado en 0) & No (uniforme) \\
Años con peso $>0$ & 1--3 por cohorte & Ninguno & Todos los pre-período \\
Concentración & Alta (1 año domina) & --- & Nula (uniforme) \\
Peso máximo típico & 0.86--1.00 & 0 & $1/T_0$ \\
\bottomrule
\end{tabular}
\end{table}

\subsection{Resumen de pesos $\hat{\omega}$}

\begin{table}[H]
\centering
\caption{Características de $\hat{\omega}$ por método}
\begin{tabular}{lccc}
\toprule
\textbf{Característica} & \textbf{SDID} & \textbf{SC} & \textbf{DiD} \\
\midrule
Optimización & Sí (regularizada) & Sí (sin regularizar) & No (uniforme) \\
Estados con peso $>0$ & $\sim$25--35 de 38 & 3--8 de 38 & Los 38 \\
Concentración & Baja-moderada & Alta & Nula \\
Peso máximo típico & 0.10--0.27 & 0.38--0.92 & 0.0263 \\
Riesgo overfitting & Bajo & Alto & Muy bajo \\
\bottomrule
\end{tabular}
\end{table}

\begin{resultado}
Los tres métodos producen ATT \textbf{idénticos a 4 decimales} entre Python y Stata,
confirma\-do por las Tablas 7--13. Las diferencias residuales de $\sim10^{-8}$ en los
pesos son acumulaciones de redondeo de punto flotante en el algoritmo Frank-Wolfe,
sin ningún efecto práctico sobre las estimaciones.

\medskip
La elección entre SDID, SC y DiD no afecta la \textit{equivalencia numérica}
entre implementaciones, pero sí afecta la \textit{estimación del ATT}:
el DiD sobreestima el efecto ($-19.48$ vs $-13.85$ SDID) porque usa controles
no comparables; el SC produce $-13.37$, muy cercano al SDID gracias a sus
pesos óptimos, pero con mayor riesgo de overfitting por su alta concentración.
\end{resultado}

% ============================================================
\newpage
\begin{thebibliography}{9}
\bibitem{arkhangelsky2021}
Arkhangelsky, D., Athey, S., Hirshberg, D.A., Imbens, G.W. \& Wager, S. (2021).
\textit{Synthetic Difference-in-Differences}.
\textit{American Economic Review}, 111(12), 4088--4118.

\bibitem{abadie2010}
Abadie, A., Diamond, A. \& Hainmueller, J. (2010).
\textit{Synthetic Control Methods for Comparative Case Studies:
Estimating the Effect of California's Tobacco Control Program}.
\textit{Journal of the American Statistical Association}, 105(490), 493--505.

\bibitem{pailanir2022}
Pailañir, D. \& Clarke, D. (2022).
\textit{SDID: Stata module to implement Synthetic Difference-in-Differences}.
Statistical Software Components S459058, Boston College.

\bibitem{baltagi2002}
Baltagi, B.H. (2002).
\textit{Econometrics} (3rd ed.). Springer.
\end{thebibliography}

\end{document}
